\section{Introduction}
\label{sec:intro}

The cache of a softmax-attention layer stores a key and a value for every token of the context
\citep{vaswani2017} and therefore grows linearly with context length. For a model shaped like
Kimi~K3.1~Pro, we estimate that the attention cache and index of a single 1M-token session occupy
32.3~GB of accelerator memory (Section~\ref{sec:serving}). Sparse attention reduces the portion of
the cache that is read at each decoding step and leaves its size unchanged \citep{nemotron35}.
Recent open hybrid models reduce the cache size by mixing softmax-attention layers with a larger
number of linear-attention layers, whose state has a fixed size. Qwen3.8-Flash-Next retains
softmax attention in 12 of its 48 layers and uses Gated DeltaNet in the remainder
\citep{qwen38fn,qwen38fncard,yang2025gdn}. Kimi~K3.1~Pro has 69 layers of Kimi Delta Attention and 24 of softmax
attention \citep{kimik31,kimilinear}, and Ling~3.1 Flash uses a 7:1 ratio \citep{ling31flash}.
GLM-5.3-Flash is the first model of the GLM series to combine linear and sparse attention, with
three linear-attention layers for each sparse-attention layer \citep{glm53flash}. The
softmax-attention layers of Qwen3.8-Flash-Next and Kimi~K3.1~Pro are likewise sparse as released,
and MiniMax-M3.1-Flash uses sparse attention in every layer, with no linear-attention layer
\citep{minimaxm31flash}.

One role of the remaining softmax-attention layers is the retrieval of earlier tokens. The
Qwen3.8-Flash-Next report states that these layers ``retain direct token-level retrieval that is
difficult for any finite-state recurrent memory to reproduce exactly'' \citep{qwen38fn}. The
GLM-5.3-Flash announcement assigns ``local dependencies'' to linear attention and the retrieval
of ``relevant global context'' to sparse attention \citep{glm53flash}. Agent
trajectories frequently repeat their own earlier text: in SWE-rebench trajectories of 16,384
tokens, a lookup that copies the continuation of the longest earlier match predicts 78.8\% of next
tokens (Section~\ref{sec:repeats}).

Two other mechanisms use the statistics of earlier text without softmax attention, static hashed
$N$-gram tables and cache language models. Static hashed $N$-gram tables, as in Engram \citep{engram2026} and the $N$-gram embedding of Qwen3.8-Flash-Next, are
parameters: they encode statistics of the training corpus and contain no information from the
current context. Cache language models \citep{kuhn1990,grave2017} interpolate the prediction of a
model with a distribution over the tokens that occurred earlier in the context. As a baseline we
use a cache of this kind that predicts the next token from counts of what followed earlier
occurrences of the current $N$-gram, and we refer to it as the n-gram cache; it operates at
inference time, outside the trained model, and has no learned memory. It shows that a lookup by exact $N$-gram match supplies recall that the
softmax-attention layers of a hybrid have not learned: after 164M training tokens, a 3:1 hybrid of
30.4M parameters reaches 0.749 nats per token, and interpolating its predictions with those of the
n-gram cache reduces this to 0.627 (Section~\ref{sec:replace}). The cache, however, is not trained
and changes only the output distribution; no layer of the model can read it, and the model keeps
its softmax-attention layers and their cache.

\begin{figure}[t]
\centering
\resizebox{\textwidth}{!}{\begin{tikzpicture}[
  font=\footnotesize,
  >={Stealth[length=1.6mm]},
  blk/.style={draw=ink!55, line width=0.3pt, minimum width=1.6cm, minimum height=0.33cm, inner sep=0pt,
              font=\scriptsize},
  gdn/.style={blk, fill=ramp1!22},
  att/.style={blk, fill=cat2!55},
  ing/.style={blk, fill=cat3!55, minimum height=0.26cm},
  tok/.style={draw=ink!55, line width=0.3pt, minimum width=0.98cm, minimum height=0.46cm, inner sep=1pt,
              font=\ttfamily\scriptsize, text height=1.25ex, text depth=0.45ex},
  hid/.style={circle, draw=ink!60, line width=0.3pt, fill=ramp1!22, minimum size=0.3cm, inner sep=0pt},
  box/.style={draw=ink!60, line width=0.4pt, rounded corners=1.5pt, inner sep=3pt, align=center},
  lab/.style={font=\scriptsize, align=center, text=ink},
]

\def\dy{0.385}
\foreach \i in {1,...,8} {
  \pgfmathparse{(\i==4 || \i==8) ? 1 : 0}
  \ifnum\pgfmathresult=1
    \node[att] (a\i) at (0, \i*\dy) {attention};
  \else
    \node[gdn] (a\i) at (0, \i*\dy) {GDN};
  \fi
}
\node[lab] at (0, 3.95) {\textbf{3:1 hybrid}};
\node[lab, below=2pt of a1] {67.5~MB};

\foreach \i in {1,...,8} {
  \pgfmathsetmacro\yy{\i*\dy + (\i>1 ? 0.33 : 0)}
  \node[gdn] (b\i) at (1.85, \yy) {GDN};
}
\node[ing] (bing) at (1.85, 1*\dy + 0.36) {\ingram{}};
\node[lab] at (1.85, 3.95) {\textbf{GDN}\\[-1pt]\textbf{+ \ingram{}}};
\node[lab, below=2pt of b1] {12.7~MB};

\foreach \i in {1,...,8} {
  \pgfmathsetmacro\yy{\i*\dy + (\i>1 ? 0.33 : 0)}
  \ifnum\i=4
    \node[att] (c\i) at (3.7, \yy) {attention};
  \else
    \node[gdn] (c\i) at (3.7, \yy) {GDN};
  \fi
}
\node[ing] (cing) at (3.7, 1*\dy + 0.36) {\ingram{}};
\node[lab] at (3.7, 3.95) {\textbf{7:1 hybrid}\\[-1pt]\textbf{+ \ingram{}}};
\node[lab, below=2pt of c1] {46.2~MB};

\node[lab, text=mute] at (1.85, -0.6) {state per session at 16,384 tokens};
\node[lab] at (1.85, 4.75) {(a) Models compared};

\begin{scope}[xshift=6.75cm]
\node[tok, draw=none] (t0) at (0.0, 0) {$\cdots$};
\node[tok, fill=cat3!18] (t1) at (0.95, 0) {tmp};
\node[tok, fill=cat3!18] (t2) at (2.0, 0) {\_};
\node[tok] (t3) at (3.05, 0) {folder};
\node[tok, draw=none] (t4) at (4.0, 0) {$\cdots$};
\node[tok, fill=cat3!18] (t5) at (4.95, 0) {tmp};
\node[tok, fill=cat3!18] (t6) at (6.0, 0) {\_};
\node[tok, densely dashed] (t7) at (7.05, 0) {folder};

\foreach \j in {1,2,3,5,6} { \node[hid] (h\j) at (t\j |- 0, 0.8) {}; \draw[ink!45, line width=0.3pt] (t\j) -- (h\j); }
\node[lab, text=mute, anchor=east, align=right] at (0.6, 0.95) {hidden state\\[-1pt]after block 1};

\draw[decorate, decoration={brace, mirror, amplitude=3pt}, ink!70] (t1.south west) ++(0,-0.06) -- ($(t2.south east)+(0,-0.06)$)
  node[midway, below=4pt, lab] {cue $c$, $N = 2$};
\draw[decorate, decoration={brace, mirror, amplitude=3pt}, ink!70] (t5.south west) ++(0,-0.06) -- ($(t6.south east)+(0,-0.06)$)
  node[midway, below=4pt, lab] {same cue $c$};
\node[lab, below=7pt of t7] {predicted};

\node[box, fill=cat3!14, minimum width=2.5cm] (slot) at (3.05, 2.9)
  {slot $S_c$ of this session\\[-1pt]{\scriptsize address: hash of $c$}};
\node[circle, draw=ink!70, inner sep=0.6pt, font=\scriptsize] (plus) at (6.0, 1.95) {$+$};
\node[box, fill=ink!6] (table) at (7.55, 1.95)
  {table row $R_c$\\[-1pt]{\scriptsize parameters}};

\draw[->, cat3!70!black, line width=0.7pt] (h3) -- (slot.south)
  node[pos=0.52, left=2pt, lab, text=cat3!50!black] {write\\$S_c \leftarrow S_c + \beta\,(v - S_c)$};
\draw[->, cat1!85!black, line width=0.7pt] (slot.east) -| (plus.north);
\draw[->, ink!70, line width=0.5pt] (table.west) -- (plus.east);
\draw[->, cat1!85!black, line width=0.7pt] (plus) -- (h6)
  node[pos=0.45, left=2pt, lab, text=cat1!70!black] {read};
\node[lab, anchor=west] at (6.2, 1.2) {$x \leftarrow x + \alpha\, W_O\,(R_c + S_c)$};
\draw[->, ink!55, line width=0.4pt, densely dashed] (h6.east) to[out=0, in=90] (t7.north);

\node[lab, text=mute] at (3.5, -1.15) {the same write and read run for each cue length $N \in \{1, 2, 4, 8, 16, 32\}$};
\node[lab] at (3.5, 4.75) {(b) One write and one read of \ingram{}};
\end{scope}
\end{tikzpicture}
}
\caption{(a) The three models of our main comparisons, each with 8 blocks: a 3:1 hybrid with softmax
attention in blocks 4 and 8, a Gated DeltaNet (GDN) model with \ingram{} after block 1, and a 7:1
hybrid with \ingram{}. The number below each model is the state that a decoder stores per session
at 16,384 tokens (Section~\ref{sec:serving}). (b) \ingram{} for one cue length, on tokens of a
held-out trajectory that names a variable \texttt{tmp\_folder}. When a token follows the cue $c$,
here the two tokens \texttt{tmp} and \texttt{\_}, the model computes a note $v$ from the hidden
state at that token and writes it into the slot $S_c$, which is addressed by the hash of $c$. When
$c$ ends at a later position, the model reads $S_c$, adds the learned table row $R_c$, and adds the
result to the hidden state through a learned gate $\alpha$ and projection $W_O$; $\beta$ is a
learned write strength. In the training trajectories this cue is most often followed by
\texttt{path} and not by \texttt{folder} (Section~\ref{sec:method-rw}).}
\label{fig:overview}
\end{figure}

We introduce In-Context N-gram Memory (\ingram{}; Figure~\ref{fig:overview}), which performs this
lookup within the model and is trained end to end. At each position, and for each of six cue lengths $N$ from 1 to 32, the
model writes a note into a slot addressed by the hash of the preceding $N$ tokens. When the same
$N$-gram ends at a later position, the model reads that slot. Because the addresses depend only on
token ids, the schedule of writes and reads is known before the forward pass. Training is
therefore parallel over the sequence, and decoding requires one read and one write per cue length
for each token, independent of context length. \ingram{} adopts the hashed addressing of Engram
and adds in-context writes. It is a write-based form of $N$-gram heads \citep{akyurek2024}, which
match earlier positions through an attention mask. We insert it once, after the first block of a
Gated DeltaNet model or hybrid.

We evaluate \ingram{} in Gated DeltaNet models and hybrids of 55.9M to 115.3M parameters and
compare them with hybrids of 30.4M to 88.4M parameters that lack it, all trained
on SWE-rebench agent trajectories with a context of 16,384 tokens. We compare models at equal
training steps on eight seeds and commit each test to the repository before the runs that it adds
(Section~\ref{sec:setup}). After 164M training tokens, a model with \ingram{} and no softmax
attention reaches a loss 0.102 nats per token lower than a parameter-matched 3:1 hybrid
(Figure~\ref{fig:confirm}). A
difference remains when the hybrid is given the same hashed table without writes (0.137 nats) and
when the n-gram cache is interpolated into both models (0.023 nats). After 1.31B training tokens,
a 7:1 hybrid with \ingram{} reaches 0.448 nats per token, compared with 0.510 for the 3:1 hybrid
and 0.508 for the 7:1 hybrid without \ingram{} (Figure~\ref{fig:training}). In a session of 131,072 tokens, \ingram{}'s state in our
models is 6.0 times smaller than the cache of one softmax-attention layer
(Figure~\ref{fig:state}). Section~\ref{sec:serving} estimates the consequence for serving: if a model shaped like
Qwen3.8-Flash-Next maintains its loss with half of its softmax-attention layers and \ingram{},
the largest context that fits on a laptop with 128~GB of memory increases from 434k to 744k
tokens.
